\documentclass[11pt]{article}
\usepackage{amsmath,amssymb}
\usepackage[margin=1in]{geometry}
\usepackage{booktabs}
\usepackage{tabularx}
\usepackage{hyperref}
\usepackage{xcolor}
\hypersetup{colorlinks=true, linkcolor=blue!50!black, urlcolor=blue!50!black}

\newcommand{\code}[1]{\texttt{#1}}
\newcommand{\QQ}{\mathbb{Q}}

\title{Does Robertz (2014) anticipate \emph{A New Method of Solving PDEs}?\\
       \large Prior-art note on \emph{Formal Algorithmic Elimination for PDEs}, LNM~2121}
\author{}
\date{25 September 2026}

\emergencystretch=2em
\begin{document}
\maketitle

\noindent
A companion note to \emph{A New Method of Solving PDEs}. The question is
whether D.~Robertz, \emph{Formal Algorithmic Elimination for PDEs}, Lecture
Notes in Mathematics~2121 (Springer, 2014) --- which the paper already cites
as \verb|\cite{RobertzLNM}|, but only for Props.~2.2.50 and~2.2.72 ---
anticipates what the paper claims is new. Page numbers below are the printed
page numbers of the book; every passage was read in the publisher's PDF text
layer (\code{\textasciitilde/library/robertz-2014/}).

\section*{Summary}

\textbf{Partly.} Robertz anticipates the \emph{strategy}: adjoin to the PDE a
system of differential equations whose solutions are exactly the functions
of a prescribed special form, run Thomas' algorithm (or a Janet basis) on the
combined system, and read off the solutions of that form. He presents this
explicitly as a generalization of separation of variables, traces it to
Olver--Rosenau (1986), and demonstrates it on the Laplace equation with
Bessel functions, on KdV, and on incompressible Navier--Stokes using the Maple
package \code{DifferentialThomas}.

He does \textbf{not} anticipate the parts that make NewMethod an algorithm
for \emph{finding} ans\"atze that work:
\begin{enumerate}
\item an ansatz carrying finitely many \emph{undetermined constants} ---
  in the argument of the unknown function ($v = v_1x+v_2y+v_3z+v_4r$) and in
  the coefficients of the ODE it satisfies --- with the output being the
  \emph{locus of constants} (consistency / membership loci) for which the
  ansatz solves the PDE;
\item the completeness theorem over constant space
  (the paper's Theorem \verb|thm:completeness|: no valid constants
  missed, no spurious ones admitted);
\item the computational route: constants kept in the coefficient field
  $\QQ(c)$, Ritt reduction modulo the Thomas cells of the ansatz, projection
  and a commutative minimal-prime decomposition (GTZ); and the ODE-extension
  encoding in which $\Psi,\Psi',\dots$ and $v$ are indeterminates tied by the
  chain rule.
\end{enumerate}

\noindent
The practical consequence is for the paper's related-work positioning, not
its results: the Conclusion currently places the method against separation
of variables and the similarity-reduction methods only, and omits the
side-condition tradition (Olver--Rosenau; Robertz \S\S3.2.5, 3.3.5) to which
it most directly belongs. See \S\ref{sec:paper}.

\section{What NewMethod claims}

From the Introduction and Conclusion of \code{NewMethod.tex}:
\begin{itemize}
\item An algorithm ``that tests a system of differential equations \dots\
  against a second such set (which I term an `ansatz'), parameterized by a
  finite number of constants \dots\ and finds'' the constant values for which
  the ansatz solves the system (Introduction, after the $J_0$ example).
\item The core algorithm: Thomas decomposition of the ansatz, Ritt reduction
  of each equation of the system modulo each component, projection onto the
  space of constants; proved to return ``exactly what they claim, neither
  missing valid constants nor admitting spurious ones''
  (\S\,Organization of the Paper; Theorem \verb|thm:completeness|).
\item Positioning (Conclusion, \code{NewMethod.tex} ll.~4126--4139): the
  technique ``is complementary to separation of variables'' and to the
  similarity-reduction methods (Lie, Clarkson--Kruskal, Bluman--Cole); its
  reductions ``are instead organized around the differential-algebraic
  structure of the ansatz \dots\ and they carry the completeness guarantee.''
\end{itemize}

\section{What Robertz already has}

\subsection{The strategy, stated in general}

Chapter~3 is about \emph{implicitization}: given a family of analytic
functions of a prescribed form, compute polynomial PDEs (and inequations)
whose analytic solutions are exactly that family. The families of \S3.2 are
linear, $\sum_i f_i(\alpha_i(z))\,g_i(z)$; those of \S3.3 are multilinear,
$p(f_1(\alpha_1(z)),\dots,f_k(\alpha_k(z)))$, where $p$ and the
\emph{inner functions $\alpha_i$ are given} and the $f_i$ are arbitrary.

The use of this for solving PDEs is stated three times:
\begin{itemize}
\item \textbf{Introduction, p.~2.} Adding a PDE system characterizing
  functions of a special form to the system under study ``may allow to
  extract explicit solutions of the prescribed type'' (citing \S3.3.5 and
  Navier--Stokes). ``This approach generalizes the well-known method of
  separation of variables.''
\item \textbf{\S3.2.5, p.~189.} If solutions of a special form are wanted,
  an implicit description of that form ``may be added to the given PDE
  system,'' and Thomas' algorithm on the combined system ``may result in PDEs
  that are easier to solve.''
\item \textbf{\S3.3.5, pp.~221 and 228.} The technique ``can be understood as
  a generalization of the method of separation of variables,'' citing
  Olver--Rosenau \cite{OR86} ``where this approach was examined from the
  symmetry point of view'' and Harrison--Estabrook \cite{HE71} for an
  exterior-differential-forms version; on p.~228: applying Thomas' algorithm
  to the combined system ``allows to decide whether there exist solutions of
  the prescribed type.''
\end{itemize}

\subsection{The worked examples}

\begin{description}
\item[Ex.~3.2.51 (p.~191): Laplace in cylindrical coordinates, Bessel
  functions.] This is the closest single example to the hydrogen
  computation. To solve Laplace's equation ``in terms of Bessel functions,''
  Robertz imposes on $u$ itself the Bessel equation in the $r$-direction,
  $r^2u_{rr} + r u_r + (r^2-\nu^2)u = 0$, adjoins it to Laplace, and computes a
  Janet basis. The consequences $u_{zz}=u$, $u_{\varphi\varphi}=-\nu^2u$ fall
  out --- ``a separation of variables'' --- giving $e^{\pm z \pm i\nu\varphi}
  J_{\pm\nu}(r)$. \\
  \emph{Relation to NewMethod:} the move ``require the solution to satisfy a
  chosen ODE, then eliminate'' is here. But the ODE is imposed on $u$ along a
  \emph{fixed coordinate} $r$ (a multiplicative separation), not on an
  unknown $\Psi$ composed with an argument $v$ to be determined; and $\nu$ is
  a symbol carried through, not a constant solved for. Nothing in the example
  would search the direction of the argument, which is what produces
  $J_0(2\sqrt{x+r})$.

\item[Ex.~3.3.49 (p.~229): KdV.] Product ansatz $u=f(t)g(x)$, encoded by the
  $2\times2$ Wronskian-type condition $u\,u_{tx}-u_t u_x=0$ (Prop.~3.3.7);
  Thomas decomposition of the combined system yields
  $u = (x+c_1)/(-6t+c_2)$ and other solutions that Maple's \code{pdsolve}
  misses.

\item[Ex.~3.3.50 (pp.~230--231): incompressible Navier--Stokes.] Each velocity
  component assumed to be a product of functions of $r$, $\varphi$, $z$
  (cylindrical), arbitrary in $t$; implicit equations adjoined; the Thomas
  decomposition is too large to finish, so simple systems are harvested one at
  a time with \code{DifferentialThomas}'s stop-and-resume feature and passed to
  \code{pdsolve}, giving an explicit solution family. \\
  \emph{Relation to NewMethod:} the same pairing (a differential-algebraic
  ansatz against Navier--Stokes, via Thomas) as the paper's
  \S\,Example: Incompressible Navier--Stokes. That section is currently
  switched off (\verb|\NSexamplefalse|, \code{NewMethod.tex} l.~126), but the
  Conclusion still names Navier--Stokes as a target (l.~4141). If the section
  is re-enabled, Ex.~3.3.50 should be cited there.

\item[Alg.~3.3.45 (p.~226): recognizing products of analytic functions.]
  Given a PDE system, decide whether its \emph{entire} solution set is of the
  form $f_1(\alpha_1)\cdots f_k(\alpha_k)$ and, if so, \emph{compute the
  linear arguments $\alpha_i$}: substitute $u=\exp(v)$, reduce to
  constant-coefficient operators, primary-decompose (citing GTZ), and read the
  $\alpha_i$ off the degree-one prime components. \\
  \emph{Relation to NewMethod:} the only place in the book where an inner
  argument is \emph{found} rather than given, and it uses GTZ to do it. But it
  applies only when the whole solution set has product form and the operators
  have constant coefficients; it does not search a parameterized ansatz for
  the sub-families that solve an arbitrary PDE.

\item[Ex.~3.2.49 (p.~190).] A Thomas decomposition shows that the
  separated-variables answer \code{pdsolve} returns for a nonlinear PDE misses
  solutions. ``Completeness'' here means \emph{completeness of a proposed
  solution family}, decided by implicitization --- a different notion from
  NewMethod's completeness over constant space.
\end{description}

\subsection{Encoding ``function of a given argument''}

Lemma~3.2.6 (p.~163): the functions of the form $f(\alpha(z))$ are the
analytic solutions of a system of first-order linear PDEs (annihilators of
$\alpha$); Remark~3.3.3 uses this to build the implicitization as an
elimination in $K\{F_1,\dots,F_k\}/Z$. NewMethod's ODE extension instead keeps
$\Psi,\Psi',\dots,\Psi^{(n)}$ and $v$ as indeterminates and defines the
$\Delta$-derivatives on them by the chain rule, so that the argument $v$ can
itself carry undetermined constants. Different encoding; same underlying
fact.

\section{What Robertz does not have}

\begin{center}
\begin{tabularx}{\linewidth}{@{}p{3.3cm}XX@{}}
\toprule
 & Robertz 2014 (\S\S3.2.5, 3.3.5) & NewMethod \\
\midrule
Inner argument & given ($\alpha_i$ fixed; found only in Alg.~3.3.45's
  constant-coefficient special case) & parameterized
  ($v=\sum v_i x_i + v_4 r$), solved for \\
Outer function & arbitrary, or a fixed ODE (Ex.~3.2.51) & satisfies an ODE
  whose coefficients are undetermined constants \\
Question answered & does the PDE have solutions of \emph{this} form, and what
  are they & for \emph{which constants} does the ansatz solve the PDE \\
Output & simple systems, then \code{pdsolve} / hand solution &
  consistency and membership loci in constant space, as minimal primes \\
Guarantee & Thomas partition of one fixed system & completeness theorem over
  the constants (no missed, no spurious $c^*$) \\
Computational route & Thomas / Janet on the combined system &
  Thomas cells of the ansatz, Ritt reduction over $\QQ(c)$, projection, GTZ \\
\bottomrule
\end{tabularx}
\end{center}

The first row is the substantive difference. With fixed arguments the method
can confirm or refute a guessed form; with parameterized arguments it
\emph{discovers} the form, which is exactly how the hydrogen computation
reaches the parabolic-coordinate solution without being told to look there.

One point that is \emph{not} new and should not be claimed: modelling
parameters as differential indeterminates with zero derivatives, ranked
lowest, is already done by Lange-Hegermann and Robertz in the control-systems
setting (\verb|\cite{LHRcontrol}|, noted in the June 2026 Thomas-completeness
session). The novelty is the projection onto constant space and the
completeness theorem, not the encoding of the constants.

\section{Suggested changes to the paper}
\label{sec:paper}

\begin{enumerate}
\item \textbf{Conclusion, related-work paragraph (ll.~4126--4139).} Add the
  side-condition / differential-constraint tradition: Olver--Rosenau
  \cite{OR86} and Robertz \verb|\cite{RobertzLNM}| \S\S3.2.5, 3.3.5. A
  framing that is accurate and keeps the claim: \emph{Adjoining to a PDE the
  differential equations characterizing functions of a prescribed form, and
  computing the consequences by Thomas' algorithm, is an established way of
  finding solutions of that form (Olver--Rosenau; Robertz). There the form is
  fixed in advance. Here the ansatz carries finitely many undetermined
  constants, the algorithm computes the locus of constants for which it solves
  the system, and the completeness theorem guarantees that locus is
  exact.}
\item \textbf{Hydrogen section.} Consider citing Robertz Ex.~3.2.51 (Laplace,
  Bessel functions via an imposed ODE) as the nearest precedent, and noting
  that the argument there is a fixed coordinate.
\item \textbf{Navier--Stokes section} (if \verb|\NSexampletrue| is restored):
  cite Ex.~3.3.50 as a previous Thomas-decomposition attack on Navier--Stokes
  with a product ansatz.
\item \textbf{The phrase ``complementary to separation of variables.''}
  Robertz presents the same strategy as a \emph{generalization} of separation
  of variables. The sentence is defensible, but a referee who knows the book
  will read it as overlooking \S3.3.5; the rewording in item~1 closes that.
\end{enumerate}

\noindent
Bibliography entries to add (verified against Robertz's reference list):

\begin{thebibliography}{9}
\bibitem{OR86} P.~J.~Olver, P.~Rosenau, The construction of special
  solutions to partial differential equations, \emph{Phys.\ Lett.\ A}
  114(3) (1986) 107--112.
\bibitem{HE71} B.~K.~Harrison, F.~B.~Estabrook, Geometric approach to
  invariance groups and solution of partial differential systems,
  \emph{J.\ Math.\ Phys.} 12 (1971) 653--666.
\end{thebibliography}

\noindent
Robertz also cites, for separation of variables in this context,
W.~Miller Jr., \emph{Symmetry and Separation of Variables} (1977);
Johnson--Reinhart--Rubel, \emph{J.\ Differential Equations} 121 (1995) 42--66;
and Gauchman--Rubel, \emph{Linear Algebra Appl.} 125 (1989) 19--63.
DOIs for \cite{OR86} and \cite{HE71} were not checked for this note.

\vfill
\noindent\small\emph{/Claude Opus 5.5 --- prepared at the author's request;
all page references checked against the publisher's PDF text layer.}

\end{document}
